\subsection{Admissible replacements with a prescribed first skeleton}

We use the finite-limit coskeleton construction of
\cite[Section 14.19, Lemmas 14.19.2 and 14.19.4]{Root_Coskeleton}.
The formulas below describe every limit and map needed here;
the admissibility comparisons are proved directly.

For $N\geq1$ define a one-truncated simplicial strict PD cochain
$S$-algebra $P(N)$ by
\[
 \begin{split}
 P(N)_0&=S,\\
 P(N)_1&=S\otimes
       \bigotimes_{w\geq N}\bigl(\Gamma_{\mathbb Z}(x_w)
                                      \otimes\Lambda_{\mathbb Z}(y_w)\bigr),
 \quad dx_w=y_w,\quad dy_w=0,\\
 |x_w|&=0,\quad |y_w|=1,\quad
 \operatorname{wt}(x_w)=\operatorname{wt}(y_w)=w,\\
 d_1x_w&=0,\quad d_0x_w=T^{[w]},\quad d_0y_w=d_1y_w=0.
 \end{split}                                                     \tag{SK.1}
\]
The degeneracy $s_0:S\to P(N)_1$ is coefficient inclusion.
Tensor products mean unions of finite tensor products; no completion
is taken. The endpoint maps in (SK.1) are the strict PD maps
specified by the displayed generator values.

\begin{theorem}[A-004: Proved; a common first-skeleton target replacement]
\label{A004}
For any pair $A,B\in\mathcal A_{\mathbb Z}$, choose
\[
                 N\geq\max\{1,\mu_A(0),\mu_B(0)\}.
 \tag{SK.2}
\]
There are explicitly constructed admissible objects $A',B'$ and
actual strict coefficient-preserving PD maps
\[
                        A'\longrightarrow A,
                 \qquad B'\longrightarrow B                     \tag{SK.3}
\]
whose first skeletons are both the specified $P(N)$, with its
displayed coefficients, differential, faces and degeneracy.
The construction preserves $q\leq w$ when that support condition
holds. It preserves termwise flatness over $\mathbb Z$.
For the object $A'$ the following bounds suffice:
\[
 \begin{split}
 \mu_{A'}(d)&=\max\{N,\mu_A(d)\},\\
 \lambda_{A'}(d)&=
       \max\{N,\lambda_A(d+2),\mu_A(d+2)\}.
 \end{split}                                                     \tag{SK.4}
\]
The analogous formulas hold for $B'$.
In the latter weight range the map $A'\to A$ induces an
isomorphism on vertical cohomology through degree $d$ in every
simplex; this is an assertion about the displayed strict map.
\end{theorem}

\begin{proof}
First, $S\to P(N)_1$ is an integral quasi-isomorphism in every
weight. A single factor in (SK.1) is the direct sum of its constants
and the identity pairs
\[
              \gamma_{r+1}(x_w)\longmapsto\gamma_r(x_w)y_w
                       \qquad(r\geq0).
 \tag{SK.5}
\]
Their contractions are integral and $S$-linear. Tensoring them
contracts every summand with a nonconstant factor; each element
uses only finitely many factors. This proves the assertion,
including weight zero. The endpoint map
\[
               P(N)_1^{w,q}\longrightarrow
                       P(N)_0^{w,q}\oplus P(N)_0^{w,q}
 \tag{SK.6}
\]
is surjective for $w\geq N$: in positive form degrees its target
is zero; in degree zero the coefficient degeneracy gives the
diagonal copy of $S^w$ and $x_w$ gives the independent pair
$(0,T^{[w]})$.

For each $w\geq N$, row exactness of $A^{w,0}$ gives
$a_w\in N_1(A^{w,0})$ with $d_0a_w=T^{[w]}$.
Thus $d_1a_w=0$. Assign $x_w\mapsto a_w$ and
$y_w\mapsto da_w$, and use the coefficient inclusion on $S$.
The free PD universal property gives a strict map
\[
                  p_A:P(N)\longrightarrow A_{\leq1}
 \tag{SK.7}
\]
of one-truncated simplicial PD DGAs. Its endpoint compatibility
follows from $d(T^{[w]})=0$ and $d(0)=0$.
The same independent choices for $B$ give $p_B$ on the very same
source $P(N)$. No compatibility between the two chosen collections
of elements is required.

Here is the explicit coskeleton used in the construction. For
one-truncated data $M=(M_0,M_1,d_0,d_1,s_0)$, put
\[
 C(M)_n=\left\{((v_i),(e_{ij}))\in
       M_0^{n+1}\times M_1^{\binom{n+1}{2}}:
       d_1e_{ij}=v_i,\ d_0e_{ij}=v_j\quad(i<j)\right\}.
 \tag{SK.8}
\]
An order-preserving map $\theta:[m]\to[n]$ selects vertices
$v_{\theta(i)}$ and edges $e_{\theta(i),\theta(j)}$; when
$\theta(i)=\theta(j)$ its edge value is $s_0v_{\theta(i)}$.
These rules compose exactly and give all faces and degeneracies.
Products, differential and PD operations are componentwise and
preserve the matching relations. The graph on these vertices is
connected, so all augmentations of a matching family agree;
they define its augmentation. Coefficients are its constant matching
families. The resulting object is the stated strict PD coskeleton.
It agrees literally with $M_0,M_1$ in simplex degrees zero and one,
under the displayed matching-coordinate identifications.

We record its exact row calculation. For the underlying
one-truncated module define
\[
 N_1M=\ker(d_1:M_1\to M_0),\qquad
 Z_1M=\ker(d_0:N_1M\to M_0).
\]
The normalization of (SK.8) is precisely
\[
                 0\longrightarrow Z_1M
                 \xrightarrow{\ \mathrm{inclusion}\ }N_1M
                 \xrightarrow{\ d_0\ }M_0,
 \tag{SK.9}
\]
in homological degrees two, one, zero.
Indeed normalized one-simplices are exactly $N_1M$.
A normalized two-simplex has vertices zero and edges $01,02$
zero, while its edge $12$ is an arbitrary element of $Z_1M$;
its remaining face is inclusion of that edge. In a normalized
simplex of degree at least three every vertex and edge belongs
to one of the zero faces $d_i$ with $i>0$, so every coordinate
is zero. This proves (SK.9) directly.
It is exact in positive degrees and has $H_0=\operatorname{coker}d_0$.

For $C(P(N))$ this $H_0$ is zero in every row of weight at least
$N$, by (SK.6). For $C(A_{\leq1})$ it is zero whenever the
corresponding row of $A$ is exact. Moreover the natural strict map
\[
                        \rho_A:A\longrightarrow C(A_{\leq1})
 \tag{SK.10}
\]
restricts a simplex to all its vertices and edges. In an exact
row it is surjective in every simplex: by A-001, the restriction
from $\Delta[n]$ to its simplicial one-skeleton is onto.
That one-skeleton has exactly the matching data (SK.8).
This proof includes every $n$ with the same row cutoff.

Define the claimed actual object and its two projections by the
pullback in strict simplicial augmented PD DGAs
\[
 A'=A\mathop{\times}_{C(A_{\leq1})} C(P(N)),
 \qquad A'\longrightarrow A,\quad A'\longrightarrow C(P(N)).
 \tag{SK.11}
\]
All its structure is the componentwise structure on pairs with
equal images. The first two simplex degrees of $\rho_A$ are
identities. For $j=0,1$ the strict identification with $P(N)_j$
has forward map $c\mapsto(p_A(c),c)$ into $A'_j$, and inverse
the second projection. We transport the structure through these
canonical strict isomorphisms in those two degrees. Thus (SK.11)
has the prescribed first skeleton, with every face and degeneracy
specified, and its map to $A$ restricts to (SK.7).

For $q\leq d$ and $w\geq\max\{N,\mu_A(d)\}$, (SK.10) is
surjective in every simplex of that row. Hence there is the
actual short exact sequence of simplicial modules
\[
 0\longrightarrow A'^{w,q}\longrightarrow
 A^{w,q}\oplus C(P(N))^{w,q}
 \xrightarrow{\ \rho_A-C(p_A)\ }C(A_{\leq1})^{w,q}
 \longrightarrow0.
 \tag{SK.12}
\]
The middle and right rows are exact by (SK.9) and the hypothesis
on $A$. Exactness of normalization and its long exact homology
sequence prove horizontal exactness of the left row, with exactly
the first bound in (SK.4).

We prove the vertical comparison without assuming either coskeleton
is a coefficient resolution. Fix $n$, and write $v=n+1$ and
$e=\binom{n+1}{2}$. Matching families give the kernel sequence
of cochain modules
\[
 0\longrightarrow C(M)_n\longrightarrow M_0^v\oplus M_1^e
 \xrightarrow{\ D_M\ }M_0^{2e},
 \qquad
 D_M((v_i),(e_{ij}))=(d_1e_{ij}-v_i,d_0e_{ij}-v_j)_{i<j}.
 \tag{SK.13}
\]
The last map is onto whenever the endpoint map $M_1\to M_0^2$
is onto, by setting all proposed vertices to zero and choosing
each edge separately.
For $M=P(N)$ this holds in weights at least $N$ by (SK.6).
For $M=A_{\leq1}$ it holds in every exact horizontal row:
the prescribed endpoints are restrictions to
$\partial\Delta[1]\subset\Delta[1]$, so A-001 applies.

Take $w$ above the second bound in (SK.4).
These surjectivities hold in all form degrees through $d+2$.
The maps $P(N)_0\to A_0$ and $P(N)_1\to A_1$ induce
isomorphisms on cohomology through $d+2$: both coefficient
comparisons have been verified, and (SK.7) fixes coefficients.
Their finite products do so as well, regardless of $v,e$.
Apply the long exact sequences to (SK.13), replacing its last
term outside this range by its actual image if necessary.
The image and the displayed last term agree through degree
$d+2$, so their cohomology agrees through degree $d+1$.
It follows that the actual map
\[
                    C(p_A)_n:C(P(N))_n\longrightarrow C(A_{\leq1})_n
 \tag{SK.14}
\]
induces isomorphisms through cohomological degree $d+1$.
In particular this comparison includes the graph's degree-one
cohomology on both sides; neither side has been declared to have
vanishing degree-one cohomology.

Finally (SK.12), now in a fixed simplex and as a vertical
cochain sequence, is short exact through degree $d+2$.
Compare it with the split sequence
\[
 0\to A_n\to A_n\oplus C(A_{\leq1})_n
                    \xrightarrow{\ \rho_A-\mathrm{id}\ }
                       C(A_{\leq1})_n\to0,
\]
using $\mathrm{id}\oplus C(p_A)_n$ on its middle terms.
For a short exact sequence in all degrees take the actual image
of $\rho_A-C(p_A)$ as its right term; its comparison is the
inclusion into $C(A_{\leq1})_n$, equal to the identity through
degree $d+2$. Its left comparison is precisely
$A'_n\to A_n$. The long exact sequences and (SK.14) prove
that this left comparison induces isomorphisms through degree
$d$. Its action on the coefficient copy is the identity.
This proves the second bound in (SK.4), uniformly in $n$.

The support condition is preserved by all displayed finite limits.
For flatness over $\mathbb Z$, the terms of $P(N)_0,P(N)_1$ are
free abelian, so the matching subgroups $C(P(N))_n$ are torsionfree.
The pullback is a subgroup of $A_n\oplus C(P(N))_n$.
When $A$ is termwise flat, that direct sum is torsionfree,
and so is its subgroup; it is flat over $\mathbb Z$.
This proves every assertion for $A'$.
The construction and argument for $B'$ use the same $P(N)$
and its independently chosen map $p_B$.
\end{proof}

\begin{corollary}[A-005: Proved; reduction to targets with a common prescribed first skeleton]
\label{A005}
The common-source assertion for all of $\mathcal A_{\mathbb Z}$
is equivalent to that assertion just for pairs whose first
skeletons are both $P(N)$ for one common integer $N\geq1$.
The same equivalence holds with support $q\leq w$, with
termwise flatness, and with both restrictions together.
The source in the restricted assertion is an object of the
corresponding full admissible category; its first skeleton is
not prescribed by this corollary.
\end{corollary}
\begin{proof}
The assertion for all pairs implies its restriction to the
displayed pairs. Conversely, for given $A,B$ apply A-004
to construct the actual arrows (SK.3). A common admissible
source $D\to A',B'$ supplied for that restricted pair gives
the strict composites $D\to A'\to A$ and $D\to B'\to B$.
Their gradings, products, PD operations and coefficient maps
are preserved by composition. Every optional restriction is
preserved by A-004 and is stipulated for $D$ in the restricted
assertion. This proves both implications with their actual maps.
\end{proof}

The target restriction is proper even up to isomorphism.
For an explicit example take the standard admissible theory $G$
and adjoin a free vertical even-to-odd identity disk on the
contractible simplicial parameter $D(1)$, in forms zero and one
and weight one. In each simplex its integral disk contraction
preserves the vertical comparison of $G$. The parameter's
simplicial contraction gives a deformation retraction of every
horizontal row onto that of $G$, by the generator formula of
AT13. Thus the resulting $F$ is admissible, flat, and supported
in $q\leq w$, with actual split strict maps $G\to F\to G$.
Its simplex-zero term contains the nonzero weight-one one-form
of that disk, whereas $P(N)_0=S$ has zero one-form group.
A form-preserving isomorphism of these simplex-zero DGAs
would have to identify those two groups, which it cannot do.
This supplies an actual object outside the restricted target
class as well as the maps relating it to a standard theory.

The corollary reduces the targets in the common-source question;
it constructs no common source for an arbitrary restricted pair.
It also makes no claim that the localization has trivial loops.
